% GENERATED FILE. Edit canonical lessons, then run npm run generate:books.
% canonical-source-hash: fnv1a64:0d294eba3ad28edd
% canonical-lessons: BN-C139-bhitu, BN-C139-shanto, BN-C139-joruri, BN-C139-soja, BN-C139-baka

\chapter{Timid, Calm, Urgent, Straight, Crooked}
\label{ch:bn-timid-calm-urgent-straight-crooked}
\hlchaptermodality{139}

\begin{chapteropening}
\textbf{By the end of this chapter:} \emph{I can say timid, calm, urgent, straight and crooked.}

Complete the last lesson of chapter 139: I can say timid, calm, urgent, straight and crooked.
\end{chapteropening}

\section[\texorpdfstring{\bn{ভীতু} (bhītu)}{ভীতু (bhītu)}]{\bn{ভীতু} (bhītu) — timid}

\label{lesson:BN-C139-bhitu}



\begin{quote}
Before the new one: say the Bengali for kind, then the Bengali for brave.
\end{quote}

\subsection*{You'll want to know: \bn{ভীতু}}
\textbf{\bn{ভীতু}} — \emph{bhītu} — \textquotedblleft{}timid\textquotedblright{}.

\begin{grammarlens}[title={What you've built}]
One more word for this chapter.
\end{grammarlens}

\subsection*{Your turn}
\begin{itemize}
\raggedright
  \item \emph{Say it:} \emph{bhītu}
  \item \emph{Say it:} \emph{bhītu}, once more
  \item \emph{Say it:} say \emph{bhītu}
  \item \emph{Recall:} say the Bengali for kind, then the Bengali for brave, then say \emph{bhītu} again
\end{itemize}

\begin{tcolorbox}[breakable,colback=teal!4,colframe=teal!35!black,title={Before you move on}]
What does \bn{ভীতু} mean? (\textquotedblleft{}Timid\textquotedblright{}.) Say it once more.
\end{tcolorbox}
\section[\texorpdfstring{\bn{শান্ত} (shāntô)}{শান্ত (shāntô)}]{\bn{শান্ত} (shāntô) — calm}

\label{lesson:BN-C139-shanto}



\begin{quote}
Before the new one: say the Bengali for brave, then the Bengali for timid.
\end{quote}

\subsection*{You'll want to know: \bn{শান্ত}}
\textbf{\bn{শান্ত}} — \emph{shāntô} — \textquotedblleft{}calm\textquotedblright{}.

\begin{grammarlens}[title={What you've built}]
One more word for this chapter.
\end{grammarlens}

\subsection*{Your turn}
\begin{itemize}
\raggedright
  \item \emph{Say it:} \emph{shāntô}
  \item \emph{Say it:} \emph{shāntô}, once more
  \item \emph{Say it:} say \emph{shāntô}
  \item \emph{Recall:} say the Bengali for brave, then the Bengali for timid, then say \emph{shāntô} again
\end{itemize}

\begin{tcolorbox}[breakable,colback=teal!4,colframe=teal!35!black,title={Before you move on}]
What does \bn{শান্ত} mean? (\textquotedblleft{}Calm\textquotedblright{}.) Say it once more.
\end{tcolorbox}
\section[\texorpdfstring{\bn{জরুরি} (jôruri)}{জরুরি (jôruri)}]{\bn{জরুরি} (jôruri) — urgent}

\label{lesson:BN-C139-joruri}



\begin{quote}
Before the new one: say the Bengali for timid, then the Bengali for calm.
\end{quote}

\subsection*{You'll want to know: \bn{জরুরি}}
\textbf{\bn{জরুরি}} — \emph{jôruri} — \textquotedblleft{}urgent\textquotedblright{}.

\begin{grammarlens}[title={What you've built}]
One more word for this chapter.
\end{grammarlens}

\subsection*{Your turn}
\begin{itemize}
\raggedright
  \item \emph{Say it:} \emph{jôruri}
  \item \emph{Say it:} \emph{jôruri}, once more
  \item \emph{Say it:} say \emph{jôruri}
  \item \emph{Recall:} say the Bengali for timid, then the Bengali for calm, then say \emph{jôruri} again
\end{itemize}

\begin{tcolorbox}[breakable,colback=teal!4,colframe=teal!35!black,title={Before you move on}]
What does \bn{জরুরি} mean? (\textquotedblleft{}Urgent\textquotedblright{}.) Say it once more.
\end{tcolorbox}
\section[\texorpdfstring{\bn{সোজা} (sojā)}{সোজা (sojā)}]{\bn{সোজা} (sojā) — straight}

\label{lesson:BN-C139-soja}



\begin{quote}
Before the new one: say the Bengali for calm, then the Bengali for urgent.
\end{quote}

\subsection*{You'll want to know: \bn{সোজা}}
\textbf{\bn{সোজা}} — \emph{sojā} — \textquotedblleft{}straight\textquotedblright{}.

\begin{grammarlens}[title={What you've built}]
One more word for this chapter.
\end{grammarlens}

\subsection*{Your turn}
\begin{itemize}
\raggedright
  \item \emph{Say it:} \emph{sojā}
  \item \emph{Say it:} \emph{sojā}, once more
  \item \emph{Say it:} say \emph{sojā}
  \item \emph{Recall:} say the Bengali for calm, then the Bengali for urgent, then say \emph{sojā} again
\end{itemize}

\begin{tcolorbox}[breakable,colback=teal!4,colframe=teal!35!black,title={Before you move on}]
What does \bn{সোজা} mean? (\textquotedblleft{}Straight\textquotedblright{}.) Say it once more.
\end{tcolorbox}
\section[\texorpdfstring{\bn{বাঁকা} (bā̃kā)}{বাঁকা (bā̃kā)}]{\bn{বাঁকা} (bā̃kā) — crooked}

\label{lesson:BN-C139-baka}



\begin{quote}
Before the new one: say the Bengali for urgent, then the Bengali for straight.
\end{quote}

\subsection*{You'll want to know: \bn{বাঁকা}}
\textbf{\bn{বাঁকা}} — \emph{bā̃kā} — \textquotedblleft{}crooked\textquotedblright{}.

\begin{grammarlens}[title={What you've built}]
One more word for this chapter.
\end{grammarlens}

\subsection*{Your turn}
\begin{itemize}
\raggedright
  \item \emph{Say it:} \emph{bā̃kā}
  \item \emph{Say it:} \emph{bā̃kā}, once more
  \item \emph{Say it:} say \emph{bā̃kā}
  \item \emph{Recall:} say the Bengali for urgent, then the Bengali for straight, then say \emph{bā̃kā} again
\end{itemize}

\begin{tcolorbox}[breakable,colback=teal!4,colframe=teal!35!black,title={Before you move on}]
What does \bn{বাঁকা} mean? (\textquotedblleft{}Crooked\textquotedblright{}.) Say it once more.
\end{tcolorbox}
